\import{nat/divisibility.tex}
\import{set/logical-graph.tex}

\section{Binomial coefficients}

Pascal's triangle is built one finite row at a time.  The row at height $n$
is a function on $\suc{n}$; no infinite row is introduced.

\begin{definition}\label{pascal_row_step}
    $\pascalRowStepOp:\Set\to\Set\to\Set
        =\fun{n,p:\Set}{
            \Apply{\graphOnOp}{
                \suc{\suc{n}},\naturals,
                \fun{k:\Set}{
                    \Apply{\logicalIfSetOp}{
                        (k=\emptyset),\naturalOne,
                        \Apply{\logicalIfSetOp}{
                            (k\in\suc{n}),
                            \apply{p}{\unions{k}}+\apply{p}{k},
                            \naturalOne}}}}}$.
\end{definition}

\begin{definition}\label{pascal_row}
    $\pascalRow{n}=\Apply{\naturalRecOp}{
        \Apply{\graphOnOp}{
            \suc{\emptyset},\naturals,
            \fun{k:\Set}{\naturalOne}},
        \pascalRowStepOp,n}$.
\end{definition}

\begin{proposition}\label{pascal_row_zero}
    $\pascalRow{\emptyset}
        =\Apply{\graphOnOp}{
            \suc{\emptyset},\naturals,
            \fun{k:\Set}{\naturalOne}}$.
\end{proposition}
\begin{proof}
    Follows by \obv{pascal_row,natural_recursion_zero}.
\end{proof}

\begin{proposition}\label{pascal_row_successor}
    Let $n$ be a natural number.
    Then $\pascalRow{\suc{n}}
        =\Apply{\pascalRowStepOp}{n,\pascalRow{n}}$.
\end{proposition}
\begin{proof}
    Follows by
        \auto{pascal_row,natural_recursion_successor}.
\end{proof}

\begin{lemma}\label{pascal_row_step_type}
    Let $n$ be a natural number.
    Let $p\in\funs{\suc{n}}{\naturals}$.
    Then $\Apply{\pascalRowStepOp}{n,p}
        \in\funs{\suc{\suc{n}}}{\naturals}$.
\end{lemma}
\begin{proof}
    Show that for every $k\in\suc{\suc{n}}$ we have
        $\Apply{\logicalIfSetOp}{
            (k=\emptyset),\naturalOne,
            \Apply{\logicalIfSetOp}{
                (k\in\suc{n}),
                \apply{p}{\unions{k}}+\apply{p}{k},
                \naturalOne}}\in\naturals$.
    \begin{subproof}
        Fix $k\in\suc{\suc{n}}$.
        \begin{byCase}
            \caseOf{$k=\emptyset$.}
                $\naturalOne$ is a natural number by
                    \auto{natural_one,omega_empty,omega_successor}.
                Follows by \obv{logical_if_set_true}.
            \caseOf{$k\neq\emptyset$.}
                \begin{byCase}
                    \caseOf{$k\in\suc{n}$.}
                        $k$ is a natural number by
                            \auto{omega_successor,
                                naturals_is_transitive_set,
                                transitiveset_iff_subseteq,subseteq}.
                        Take a natural number $j$ such that $k=\suc{j}$
                            by \obv{natural_zero_or_successor}.
                        $j$ is \in-transitive by
                            \obv{natural_number_is_ordinal,
                                ordinal_is_transitiveset}.
                        Hence $\unions{k}=j$ by
                            \auto{transitiveset_iff_unions_suc}.
                        $j\in k$ by \obv{suc_intro_self}.
                        Hence $j\in\suc{n}$ by
                            \obv{omega_successor,
                                natural_number_is_ordinal,
                                ordinal_transitivity}.
                        $p(j),p(k)\in\naturals$ by
                            \obv{funs_type_apply}.
                        Hence $p(j)+p(k)$ is a natural number by
                            \obv{num_addition_on_naturals}.
                        Follows by
                            \auto{logical_if_set_false,
                                logical_if_set_true}.
                    \caseOf{$k\notin\suc{n}$.}
                        $\naturalOne$ is a natural number by
                            \auto{natural_one,omega_empty,
                                omega_successor}.
                        Follows by
                            \auto{logical_if_set_false}.
                \end{byCase}
        \end{byCase}
    \end{subproof}
    Follows by
        \auto{pascal_row_step,logical_graph_on_function}.
\end{proof}

\begin{lemma}\label{pascal_row_step_first}
    Let $n$ be a natural number and $p$ be a set.
    Then $\apply{\Apply{\pascalRowStepOp}{n,p}}{\emptyset}
        =\naturalOne$.
\end{lemma}
\begin{proof}
    $\suc{\suc{n}}$ is a natural number and
        $\suc{\suc{n}}\neq\emptyset$ by
        \obv{omega_successor,suc_neq_emptyset}.
    Hence $\emptyset\in\suc{\suc{n}}$ by
        \obv{natural_nonzero_contains_zero}.
    Follows by
        \auto{pascal_row_step,logical_graph_on_apply,
            logical_if_set_true}.
\end{proof}

\begin{lemma}\label{pascal_row_step_last}
    Let $n$ be a natural number and $p$ be a set.
    Then $\apply{\Apply{\pascalRowStepOp}{n,p}}{\suc{n}}
        =\naturalOne$.
\end{lemma}
\begin{proof}
    $\suc{n}\in\suc{\suc{n}}$ by \obv{suc_intro_self}.
    $\suc{n}\neq\emptyset$ by \obv{suc_neq_emptyset}.
    $\suc{n}\notin\suc{n}$ by \obv{in_irrefl}.
    Follows by
        \auto{pascal_row_step,logical_graph_on_apply,
            logical_if_set_false}.
\end{proof}

\begin{lemma}\label{pascal_row_step_interior}
    Let $n$ be a natural number.
    Let $p\in\funs{\suc{n}}{\naturals}$.
    Let $k\in n$.
    Then
        $\apply{\Apply{\pascalRowStepOp}{n,p}}{\suc{k}}
            =p(k)+p(\suc{k})$.
\end{lemma}
\begin{proof}
    $k$ is a natural number by
        \obv{naturals_is_transitive_set,
            transitiveset_iff_subseteq,subseteq}.
    $\suc{k}$ is a natural number by \obv{omega_successor}.
    Hence $\suc{k}$ and $n$ are ordinals by
        \obv{natural_number_is_ordinal}.
    Hence $\suc{k}\subseteq n$ by \obv{ordinal_suc_subseteq}.
    Hence $\suc{k}\in n$ or $\suc{k}=n$ by
        \obv{ordinal_subseteq_implies_elem_or_eq}.
    Hence $\suc{k}\in\suc{n}$ by \auto{suc_iff}.
    Hence $\suc{k}\in\suc{\suc{n}}$ by
        \auto{suc_iff}.
    $k$ is \in-transitive by
        \obv{natural_number_is_ordinal,ordinal_is_transitiveset}.
    Hence $\unions{\suc{k}}=k$ by
        \obv{transitiveset_iff_unions_suc}.
    $\suc{k}\neq\emptyset$ by \obv{suc_neq_emptyset}.
    Follows by
        \auto{pascal_row_step,logical_graph_on_apply,
            logical_if_set_false,logical_if_set_true}.
\end{proof}

\begin{theorem}\label{pascal_row_type}
    Let $n$ be a natural number.
    Then $\pascalRow{n}\in\funs{\suc{n}}{\naturals}$.
\end{theorem}
\begin{proof}
    Let $C=\{m\in\naturals\mid
        \pascalRow{m}\in\funs{\suc{m}}{\naturals}\}$.
    $\naturalOne$ is a natural number by
        \auto{natural_one,omega_empty,omega_successor}.
    $\pascalRow{\emptyset}\in
        \funs{\suc{\emptyset}}{\naturals}$ by
        \auto{pascal_row_zero,logical_graph_on_function}.
    Hence $\emptyset\in C$ by \obv{omega_empty}.
    Show that for every $m\in C$ we have $\suc{m}\in C$.
    \begin{subproof}
        Fix $m\in C$.
        Hence $m$ is a natural number and
            $\pascalRow{m}\in\funs{\suc{m}}{\naturals}$.
        $\Apply{\pascalRowStepOp}{m,\pascalRow{m}}
            \in\funs{\suc{\suc{m}}}{\naturals}$ by
            \obv{pascal_row_step_type}.
        Hence $\pascalRow{\suc{m}}
            \in\funs{\suc{\suc{m}}}{\naturals}$ by
            \obv{pascal_row_successor}.
        $\suc{m}$ is a natural number by \obv{omega_successor}.
    \end{subproof}
    Hence $C$ is an inductive set by \obv{inductive_set}.
    Hence every natural number is an element of $C$ by
        \obv{omega_minimal}.
    Thus $n\in C$.
\end{proof}

\begin{proposition}\label{pascal_row_first}
    Let $n$ be a natural number.
    Then $\apply{\pascalRow{n}}{\emptyset}=\naturalOne$.
\end{proposition}
\begin{proof}
    Let $C=\{m\in\naturals\mid
        \apply{\pascalRow{m}}{\emptyset}=\naturalOne\}$.
    $\emptyset\in\suc{\emptyset}$ by \obv{suc_intro_self}.
    $\apply{\pascalRow{\emptyset}}{\emptyset}=\naturalOne$ by
        \auto{pascal_row_zero,logical_graph_on_apply}.
    Hence $\emptyset\in C$ by \obv{omega_empty}.
    Show that for every $m\in C$ we have $\suc{m}\in C$.
    \begin{subproof}
        Fix $m\in C$.
        $\pascalRow{\suc{m}}
            =\Apply{\pascalRowStepOp}{m,\pascalRow{m}}$ by
            \obv{pascal_row_successor}.
        Hence $\apply{\pascalRow{\suc{m}}}{\emptyset}=\naturalOne$ by
            \auto{pascal_row_step_first}.
        $\suc{m}$ is a natural number by \obv{omega_successor}.
    \end{subproof}
    Hence $C$ is an inductive set by \obv{inductive_set}.
    Hence every natural number is an element of $C$ by
        \obv{omega_minimal}.
    Thus $n\in C$.
\end{proof}

\begin{proposition}\label{pascal_row_last}
    Let $n$ be a natural number.
    Then $\apply{\pascalRow{n}}{n}=\naturalOne$.
\end{proposition}
\begin{proof}
    Let $C=\{m\in\naturals\mid
        \apply{\pascalRow{m}}{m}=\naturalOne\}$.
    $\emptyset\in\suc{\emptyset}$ by \obv{suc_intro_self}.
    $\apply{\pascalRow{\emptyset}}{\emptyset}=\naturalOne$ by
        \auto{pascal_row_zero,logical_graph_on_apply}.
    Hence $\emptyset\in C$ by \obv{omega_empty}.
    Show that for every $m\in C$ we have $\suc{m}\in C$.
    \begin{subproof}
        Fix $m\in C$.
        $\pascalRow{\suc{m}}
            =\Apply{\pascalRowStepOp}{m,\pascalRow{m}}$ by
            \obv{pascal_row_successor}.
        Hence $\apply{\pascalRow{\suc{m}}}{\suc{m}}=\naturalOne$ by
            \auto{pascal_row_step_last}.
        $\suc{m}$ is a natural number by \obv{omega_successor}.
    \end{subproof}
    Hence $C$ is an inductive set by \obv{inductive_set}.
    Hence every natural number is an element of $C$ by
        \obv{omega_minimal}.
    Thus $n\in C$.
\end{proof}

\begin{proposition}\label{pascal_row_interior}
    Let $n$ be a natural number and $k\in n$.
    Then
        $\apply{\pascalRow{\suc{n}}}{\suc{k}}
            =\apply{\pascalRow{n}}{k}
                +\apply{\pascalRow{n}}{\suc{k}}$.
\end{proposition}
\begin{proof}
    $\pascalRow{n}\in\funs{\suc{n}}{\naturals}$ by
        \obv{pascal_row_type}.
    $\pascalRow{\suc{n}}
        =\Apply{\pascalRowStepOp}{n,\pascalRow{n}}$ by
        \obv{pascal_row_successor}.
    Follows by \auto{pascal_row_step_interior}.
\end{proof}

\subsection{Total binomial coefficients}

\begin{definition}\label{natural_choose}
    $\naturalChoose{n}{k}=\Apply{\logicalIfSetOp}{
        (n\in\naturals\land k\in\naturals\land k\in\suc{n}),
        \apply{\pascalRow{n}}{k},\emptyset}$.
\end{definition}

\begin{proposition}\label{natural_choose_in_range}
    Let $n,k$ be natural numbers and $k\in\suc{n}$.
    Then $\naturalChoose{n}{k}=\apply{\pascalRow{n}}{k}$.
\end{proposition}
\begin{proof}
    Follows by \obv{natural_choose,logical_if_set_true}.
\end{proof}

\begin{proposition}\label{natural_choose_off_domain}
    Suppose $n\notin\naturals$ or $k\notin\naturals$ or
        $k\notin\suc{n}$.
    Then $\naturalChoose{n}{k}=\emptyset$.
\end{proposition}
\begin{proof}
    Follows by \obv{natural_choose,logical_if_set_false}.
\end{proof}

\begin{proposition}\label{natural_choose_type}
    $\naturalChoose{n}{k}$ is a natural number.
\end{proposition}
\begin{proof}
    \begin{byCase}
        \caseOf{$n\in\naturals\land k\in\naturals\land
                k\in\suc{n}$.}
            $\pascalRow{n}\in\funs{\suc{n}}{\naturals}$ by
                \obv{pascal_row_type}.
            Hence $\apply{\pascalRow{n}}{k}\in\naturals$ by
                \obv{funs_type_apply}.
            Follows by \obv{natural_choose_in_range}.
        \caseOf{$\lnot(n\in\naturals\land k\in\naturals\land
                k\in\suc{n})$.}
            $\naturalChoose{n}{k}=\emptyset$ by
                \obv{natural_choose,logical_if_set_false}.
            Follows by \obv{omega_empty}.
    \end{byCase}
\end{proof}

\begin{proposition}\label{natural_choose_first}
    Let $n$ be a natural number.
    Then $\naturalChoose{n}{\emptyset}=\naturalOne$.
\end{proposition}
\begin{proof}
    $\emptyset$ is a natural number by \obv{omega_empty}.
    $\suc{n}$ is a natural number and $\suc{n}\neq\emptyset$ by
        \obv{omega_successor,suc_neq_emptyset}.
    Hence $\emptyset\in\suc{n}$ by
        \obv{natural_nonzero_contains_zero}.
    $\naturalChoose{n}{\emptyset}
        =\apply{\pascalRow{n}}{\emptyset}$ by
        \obv{natural_choose_in_range}.
    $\apply{\pascalRow{n}}{\emptyset}=\naturalOne$ by
        \obv{pascal_row_first}.
\end{proof}

\begin{proposition}\label{natural_choose_last}
    Let $n$ be a natural number.
    Then $\naturalChoose{n}{n}=\naturalOne$.
\end{proposition}
\begin{proof}
    $n\in\suc{n}$ by \obv{suc_intro_self}.
    $\naturalChoose{n}{n}=\apply{\pascalRow{n}}{n}$ by
        \obv{natural_choose_in_range}.
    $\apply{\pascalRow{n}}{n}=\naturalOne$ by
        \obv{pascal_row_last}.
\end{proof}

\begin{theorem}\label{natural_choose_pascal}
    Let $n,k$ be natural numbers.
    Then
        $\naturalChoose{\suc{n}}{\suc{k}}
            =\naturalChoose{n}{k}+\naturalChoose{n}{\suc{k}}$.
\end{theorem}
\begin{proof}
    $n$ and $k$ are ordinals by \obv{natural_number_is_ordinal}.
    \begin{byCase}
        \caseOf{$k\in n$.}
            $\suc{k}$ is a natural number by \obv{omega_successor}.
            $\suc{k}$ and $n$ are ordinals by
                \obv{natural_number_is_ordinal}.
            $\suc{k}\subseteq n$ by \obv{ordinal_suc_subseteq}.
            Hence $\suc{k}\in n$ or $\suc{k}=n$ by
                \obv{ordinal_subseteq_implies_elem_or_eq}.
            Hence $k,\suc{k}\in\suc{n}$ by \auto{suc_iff}.
            Hence $\suc{k}\in\suc{\suc{n}}$ by \auto{suc_iff}.
            $\naturalChoose{\suc{n}}{\suc{k}}
                =\apply{\pascalRow{\suc{n}}}{\suc{k}}$ and
                $\naturalChoose{n}{k}=\apply{\pascalRow{n}}{k}$ and
                $\naturalChoose{n}{\suc{k}}
                    =\apply{\pascalRow{n}}{\suc{k}}$ by
                \obv{omega_successor,natural_choose_in_range}.
            $\apply{\pascalRow{\suc{n}}}{\suc{k}}
                =\apply{\pascalRow{n}}{k}
                    +\apply{\pascalRow{n}}{\suc{k}}$ by
                \obv{pascal_row_interior}.

        \caseOf{$k\notin n$.}
            \begin{byCase}
                \caseOf{$k=n$.}
                    $\suc{n}$ is a natural number by
                        \obv{omega_successor}.
                    $\suc{n}\notin\suc{n}$ by \obv{in_irrefl}.
                    $\naturalChoose{\suc{n}}{\suc{k}}=\naturalOne$
                        by \auto{natural_choose_last}.
                    $\naturalChoose{n}{k}=\naturalOne$ by
                        \auto{natural_choose_last}.
                    $\naturalChoose{n}{\suc{k}}=\emptyset$ by
                        \auto{natural_choose_off_domain}.
                    $\naturalOne+\emptyset=\naturalOne$ by
                        \auto{natural_one,omega_empty,
                            omega_successor,num_addition_zero_right}.

                \caseOf{$k\neq n$.}
                    Hence $n\in k$ by
                        \obv{ordinal_in_trichotomy}.
                    Show that $k\notin\suc{n}$.
                    \begin{subproof}
                        Suppose $k\in\suc{n}$.
                        Then $k=n$ or $k\in n$ by \obv{suc_elim}.
                        Contradiction.
                    \end{subproof}
                    Show that $\suc{k}\notin\suc{n}$.
                    \begin{subproof}
                        Suppose $\suc{k}\in\suc{n}$.
                        Then $\suc{k}=n$ or $\suc{k}\in n$ by
                            \obv{suc_elim}.
                        $k\in\suc{k}$ by \obv{suc_intro_self}.
                        Hence $k\in n$ by
                            \obv{natural_number_is_ordinal,
                                ordinal_transitivity}.
                        Contradiction by \obv{in_asymmetric}.
                    \end{subproof}
                    Show that $\suc{k}\notin\suc{\suc{n}}$.
                    \begin{subproof}
                        Suppose $\suc{k}\in\suc{\suc{n}}$.
                        Then $\suc{k}=\suc{n}$ or
                            $\suc{k}\in\suc{n}$ by \obv{suc_elim}.
                        Hence $k=n$ or $\suc{k}\in\suc{n}$ by
                            \obv{suc_injective}.
                        Contradiction.
                    \end{subproof}
                    $\naturalChoose{\suc{n}}{\suc{k}}=\emptyset$ and
                        $\naturalChoose{n}{k}=\emptyset$ and
                        $\naturalChoose{n}{\suc{k}}=\emptyset$ by
                        \obv{natural_choose_off_domain}.
                    $\emptyset+\emptyset=\emptyset$ by
                        \obv{omega_empty,num_addition_zero_right}.
            \end{byCase}
    \end{byCase}
\end{proof}

\begin{proposition}\label{natural_choose_one}
    Let $n$ be a natural number.
    Then $\naturalChoose{\suc{n}}{\naturalOne}=\suc{n}$.
\end{proposition}
\begin{proof}
    Let $C=\{m\in\naturals\mid
        \naturalChoose{\suc{m}}{\naturalOne}=\suc{m}\}$.
    $\naturalOne=\suc{\emptyset}$ by \obv{natural_one}.
    $\suc{\emptyset}$ is a natural number by
        \obv{omega_empty,omega_successor}.
    $\naturalChoose{\suc{\emptyset}}{\suc{\emptyset}}
        =\naturalOne$ by \obv{natural_choose_last}.
    Hence $\naturalChoose{\suc{\emptyset}}{\naturalOne}
        =\naturalOne$.
    Hence $\naturalChoose{\suc{\emptyset}}{\naturalOne}
        =\suc{\emptyset}$.
    Thus $\emptyset\in C$ by \obv{omega_empty}.
    Show that for every $m\in C$ we have $\suc{m}\in C$.
    \begin{subproof}
        Fix $m\in C$.
        Hence $m$ is a natural number and
            $\naturalChoose{\suc{m}}{\naturalOne}=\suc{m}$.
        $\naturalChoose{\suc{\suc{m}}}{\suc{\emptyset}}
            =\naturalChoose{\suc{m}}{\emptyset}
                +\naturalChoose{\suc{m}}{\suc{\emptyset}}$ by
            \obv{natural_choose_pascal,omega_empty,omega_successor}.
        $\naturalChoose{\suc{m}}{\emptyset}=\naturalOne$ by
            \obv{natural_choose_first,omega_successor}.
        Hence $\naturalChoose{\suc{\suc{m}}}{\naturalOne}
            =\naturalOne+\suc{m}$.
        $\naturalOne+\suc{m}=\suc{\suc{m}}$ by
            \auto{natural_one,num_addition_successor_left,
                num_addition_zero_left,omega_successor}.
        Hence $\naturalChoose{\suc{\suc{m}}}{\naturalOne}
            =\suc{\suc{m}}$.
        $\suc{m}$ is a natural number by \obv{omega_successor}.
    \end{subproof}
    Hence $C$ is an inductive set by \obv{inductive_set}.
    Hence every natural number is an element of $C$ by
        \obv{omega_minimal}.
    Thus $n\in C$.
\end{proof}

\begin{theorem}\label{natural_choose_weighted}
    Let $m,k$ be natural numbers and $k\in\suc{m}$.
    Then
        $\naturalMul{\suc{k}}{\naturalChoose{\suc{m}}{\suc{k}}}
            =\naturalMul{\suc{m}}{\naturalChoose{m}{k}}$.
\end{theorem}
\begin{proof}
    Let $C=\{r\in\naturals\mid
        \text{for every $i\in\suc{r}$ we have
            $\naturalMul{\suc{i}}{
                \naturalChoose{\suc{r}}{\suc{i}}}
                =\naturalMul{\suc{r}}{\naturalChoose{r}{i}}$}\}$.

    Show that $\emptyset\in C$.
    \begin{subproof}
        Show that for every $i\in\suc{\emptyset}$ we have
            $\naturalMul{\suc{i}}{
                \naturalChoose{\suc{\emptyset}}{\suc{i}}}
                =\naturalMul{\suc{\emptyset}}{
                    \naturalChoose{\emptyset}{i}}$.
        \begin{subproof}
            Fix $i\in\suc{\emptyset}$.
            Hence $i=\emptyset$ by \auto{suc_iff,emptyset}.
            $\suc{\emptyset}$ is a natural number by
                \obv{omega_empty,omega_successor}.
            $\naturalChoose{\suc{\emptyset}}{\suc{\emptyset}}
                =\naturalOne$ and
                $\naturalChoose{\emptyset}{\emptyset}=\naturalOne$ by
                \obv{natural_choose_last,omega_empty}.
            $\naturalOne=\suc{\emptyset}$ by \obv{natural_one}.
            $\naturalMul{\naturalOne}{\naturalOne}=\naturalOne$ by
                \auto{natural_multiplication_one_left}.
        \end{subproof}
        Follows by \obv{omega_empty}.
    \end{subproof}

    Show that for every $r\in C$ we have $\suc{r}\in C$.
    \begin{subproof}
        Fix $r\in C$.
        Hence $r$ is a natural number and for every $i\in\suc{r}$ we have
            $\naturalMul{\suc{i}}{
                \naturalChoose{\suc{r}}{\suc{i}}}
                =\naturalMul{\suc{r}}{\naturalChoose{r}{i}}$.
        Show that for every $i\in\suc{\suc{r}}$ we have
            $\naturalMul{\suc{i}}{
                \naturalChoose{\suc{\suc{r}}}{\suc{i}}}
                =\naturalMul{\suc{\suc{r}}}{
                    \naturalChoose{\suc{r}}{i}}$.
        \begin{subproof}
            Fix $i\in\suc{\suc{r}}$.
            $i$ is a natural number by
                \auto{omega_successor,naturals_is_transitive_set,
                    transitiveset_iff_subseteq,subseteq}.
            \begin{byCase}
                \caseOf{$i=\suc{r}$.}
                    $\suc{r}$ and $\suc{\suc{r}}$ are natural numbers by
                        \obv{omega_successor}.
                    $\naturalChoose{\suc{\suc{r}}}{\suc{i}}
                        =\naturalOne$ and
                        $\naturalChoose{\suc{r}}{i}=\naturalOne$ by
                        \auto{natural_choose_last}.
                    Follows by
                        \obv{natural_one,
                            natural_multiplication_one_right}.

                \caseOf{$i\neq\suc{r}$.}
                    Hence $i\in\suc{r}$ by \obv{suc_elim}.
                    \begin{byCase}
                        \caseOf{$i=\emptyset$.}
                            $\suc{r}$ and $\suc{\suc{r}}$ are natural
                                numbers by \obv{omega_successor}.
                            $\naturalChoose{\suc{\suc{r}}}{\naturalOne}
                                =\suc{\suc{r}}$ by
                                \obv{natural_choose_one}.
                            $\naturalChoose{\suc{r}}{\emptyset}
                                =\naturalOne$ by
                                \obv{natural_choose_first}.
                            $\naturalOne=\suc{\emptyset}$ by
                                \obv{natural_one}.
                            Follows by
                                \auto{natural_multiplication_one_left,
                                    natural_multiplication_one_right}.

                        \caseOf{$i\neq\emptyset$.}
                            Take a natural number $j$ such that
                                $i=\suc{j}$ by
                                \obv{natural_zero_or_successor}.
                            $\suc{r}$ is an ordinal by
                                \obv{omega_successor,
                                    natural_number_is_ordinal}.
                            $j\in i$ by \obv{suc_intro_self}.
                            Hence $j\in\suc{r}$ by
                                \obv{ordinal_transitivity}.

                            Let $A=\naturalChoose{\suc{r}}{i}$.
                            Let $B=\naturalChoose{\suc{r}}{\suc{i}}$.
                            Let $D=\naturalChoose{r}{j}$.
                            Let $E=\naturalChoose{r}{i}$.
                            $\suc{i}$ and $\suc{r}$ are natural numbers by
                                \obv{omega_successor}.
                            $A,B,D,E\in\naturals$ by
                                \auto{natural_choose_type}.
                            $\naturalChoose{\suc{\suc{r}}}{\suc{i}}
                                =\naturalChoose{\suc{r}}{i}
                                    +\naturalChoose{\suc{r}}{\suc{i}}$ by
                                \obv{natural_choose_pascal,
                                    omega_successor}.
                            Hence $\naturalChoose{\suc{\suc{r}}}{\suc{i}}
                                =A+B$.
                            $\naturalChoose{\suc{r}}{\suc{j}}
                                =\naturalChoose{r}{j}
                                    +\naturalChoose{r}{\suc{j}}$ by
                                \obv{natural_choose_pascal}.
                            Hence $A=D+E$.
                            $\naturalMul{i}{A}
                                =\naturalMul{\suc{r}}{D}$ and
                                $\naturalMul{\suc{i}}{B}
                                    =\naturalMul{\suc{r}}{E}$ by
                                \auto.

                            $\naturalMul{\suc{i}}{A+B}
                                =\naturalMul{\suc{i}}{A}
                                    +\naturalMul{\suc{i}}{B}$ by
                                \obv{natural_multiplication_distributive_left}.
                            $\naturalMul{\suc{i}}{A}
                                =\naturalMul{i}{A}+A$ by
                                \obv{natural_multiplication_successor_left}.
                            Hence $\naturalMul{\suc{i}}{A+B}
                                =(\naturalMul{\suc{r}}{D}+A)
                                    +\naturalMul{\suc{r}}{E}$.
                            $(\naturalMul{\suc{r}}{D}+A)
                                    +\naturalMul{\suc{r}}{E}
                                =(\naturalMul{\suc{r}}{D}
                                    +\naturalMul{\suc{r}}{E})+A$ by
                                \auto{natural_addition_associative,
                                    natural_addition_commutative,
                                    natural_multiplication_closed}.
                            $\naturalMul{\suc{r}}{D}
                                    +\naturalMul{\suc{r}}{E}
                                =\naturalMul{\suc{r}}{D+E}$ by
                                \auto{natural_multiplication_distributive_left}.
                            Hence $\naturalMul{\suc{i}}{A+B}
                                =\naturalMul{\suc{r}}{A}+A$.
                            $\naturalMul{\suc{\suc{r}}}{A}
                                =\naturalMul{\suc{r}}{A}+A$ by
                                \obv{natural_multiplication_successor_left}.
                            Hence $\naturalMul{\suc{i}}{A+B}
                                =\naturalMul{\suc{\suc{r}}}{A}$.
                    \end{byCase}
            \end{byCase}
        \end{subproof}
        $\suc{r}$ is a natural number by \obv{omega_successor}.
    \end{subproof}
    Hence $C$ is an inductive set by \obv{inductive_set}.
    Hence every natural number is an element of $C$ by
        \obv{omega_minimal}.
    Hence $m\in C$.
\end{proof}

\begin{theorem}\label{prime_divides_interior_natural_choose}
    Let $p,k$ be natural numbers.
    Suppose $p$ is a prime natural number and
        $\emptyset\in k$ and $k\in p$.
    Then $p$ divides $\naturalChoose{p}{k}$ as natural numbers.
\end{theorem}
\begin{proof}
    $p\neq\emptyset$ by \obv{emptyset}.
    Take a natural number $m$ such that $p=\suc{m}$ by
        \obv{natural_zero_or_successor}.
    $k\neq\emptyset$ by \obv{emptyset}.
    Take a natural number $j$ such that $k=\suc{j}$ by
        \obv{natural_zero_or_successor}.
    $p$ is an ordinal by \obv{natural_number_is_ordinal}.
    $j\in k$ by \obv{suc_intro_self}.
    Hence $j\in p$ by \obv{ordinal_transitivity}.
    Hence $j\in\suc{m}$.

    $\naturalMul{\suc{j}}{
            \naturalChoose{\suc{m}}{\suc{j}}}
        =\naturalMul{\suc{m}}{\naturalChoose{m}{j}}$ by
        \obv{natural_choose_weighted}.
    Hence $\naturalMul{k}{\naturalChoose{p}{k}}
        =\naturalMul{p}{\naturalChoose{m}{j}}$.

    $\naturalChoose{m}{j}$ is a natural number by
        \obv{natural_choose_type}.
    $p$ divides $p$ as natural numbers by
        \obv{natural_divides_itself}.
    Hence $p$ divides
        $\naturalMul{p}{\naturalChoose{m}{j}}$ as natural numbers by
        \obv{divisor_divides_multiple}.
    Hence $p$ divides
        $\naturalMul{k}{\naturalChoose{p}{k}}$ as natural numbers.
    Hence $p$ divides $k$ as natural numbers or
        $p$ divides $\naturalChoose{p}{k}$ as natural numbers by
        \obv{prime_divides_natural_product,natural_choose_type}.
    $p$ does not divide $k$ as natural numbers by
        \obv{larger_natural_does_not_divide_positive_smaller}.
\end{proof}
